\section*{Exercises about Invariant Subspaces}

\exercise{1}
\begin{xrcs}
  Suppose $T \in \linmap(V)$ and $U$ is a subspace of $V$. Then
  \begin{enumerate}
    \item $U \subseteq \mynull T \implies$ $U$ is invariant under $T$.
    \item $\myrange T \subseteq U \implies$ $U$ is invariant under $T$.
  \end{enumerate}

  \begin{xprf}
    { }(a). { }If $v \in U$, then $Tv = 0 \in U$, since $U$ is a subspace and therefore $0 \in U$. \\
    (b). { }Let $w \in U\subseteq V$. $\implies$ $Tw \in \myrange T$, and since $\myrange T \subseteq U$, $Tw \in U$.
  \end{xprf}

\end{xrcs}

\exercise{2}
\begin{xrcs}
  Suppose that $T \in \linmap (V)$ and $V_1, \ddd, V_m$ are subspaces of $V$ invariant under $T$.
  Prove that $V_1 + \cdots + V_m$ is invariant under $T$.

  \begin{xprf}
    Let $w \in V_1 + \cdots + V_m$ \st
    $
    w = v_1 + \cdots + v_m \mytext{where} v_1 \in V_1, \ddd, v_m \in V_m.
    $
    Then $Tw = Tv_1 + \cdots + Tv_m$ such that $Tv_1 \in V_1, \ddd, Tv_m \in V_m$ because $V_1, \ddd, V_m$ are invariant under $T$. Therefore we have $Tw \in V_1 + \cdots + V_m$. So $V_1 + \cdots + V_m$ is invariant under $T$ as well.
  \end{xprf}
\end{xrcs}

\exercise{3}
\begin{xrcs}
  Suppose $T \in \linmap(V)$. Prove that the intersection of every collection of subspaces of $V$ invariant under $T$ is invariant under $T$.

  \begin{xprf}
    Let $w \in V_1 \cap \cdots \cap V_m \subseteq V$ where each $V_i$ is invariant under $T$. Then $w \in V_1, \ddd, w \in V_m$, and thus, $Tw \in V_1, \ddd, Tw \in V_m$. Hence, $Tw \in V_1 \cap \cdots \cap V_m$.
  \end{xprf}
\end{xrcs}

\exercise{12}
\begin{xrcs}
  Suppose $V = U \oplus W$, where $U$ and $W$ are nonzero subspaces of $V$. Define $P \in \linmap(V)$ by $P(u+w) = u$ for each $u \in U$ and each $w \in W$. Find all eigenvalues and eigenvectors of $P$.

  \begin{xsol}
    We have an eigenvalue $\lambda_1 = 1$ for all vectors $u \in U$, because $P$ acts as the identity on $U$. We also have an eigenvalue $\lambda_0 = 0$ for all $w \in W$, because $P$ acts as the zero operator there.
  \end{xsol}
\end{xrcs}

\exercise{15}
\begin{xrcs}
  Suppose $V$ is finite-dimensional, $T \in \linmap (V)$, and $\lambda \in \myF$. Show that $\lambda$ is an eigenvalue of $T$ $\iff$ $\lambda$ is an eigenvalue of the dual operator $\dual{T} \in \linmap(V')$.

  \begin{xsol}
    The definition of $\dual{T}$ states that
    \begin{equation}
      \dual{T}(\varphi) = \varphi \circ T \quad \forall \varphi \in \dual{V}.
    \end{equation}

    Hence, being an eigenvalue $\lambda$ and an eigenvector $\varphi \in \dual{V} = \linmap(V, \myF)$ of $\dual{T}$ means:
    \begin{equation}
      \dual{T}(\varphi) = \lambda \varphi.
    \end{equation}

    $\lambda$ is an eigenvalue of $T$ $\iff$ $T-\lambda I$ is not injective $\iff$ $\dual{(T-\lambda I)}$ is not surjective $\iff$ $\dual{T} - \lambda I$ is not invertible $\iff$ $\lambda$ is an eigenvalue of $\dual{T}$.
  \end{xsol}
\end{xrcs}

\begin{ainote}[Do not try to build the eigenvector of $\dual{T}$ by hand]
  A tempting route is to start from an eigenvector $u$ of $T$, extend $u$ to a basis $u, v_2, \ddd, v_n$ of $V$, and define $\varphi \in \dual{V}$ by
  \begin{equation}
    \varphi(u) := 1, \qquad \varphi(v_i) := 0 \quad (i = 2, \ddd, n),
  \end{equation}
  hoping that $\varphi$ is an eigenvector of $\dual{T}$. It is not, in general. Checking $\dual{T}\varphi = \lambda \varphi$ means checking it on \emph{every} basis vector, and while
  \begin{equation}
    (\dual{T}\varphi)(u) = \varphi(Tu) = \varphi(\lambda u) = \lambda = \lambda \varphi(u)
  \end{equation}
  works out, the remaining conditions $(\dual{T}\varphi)(v_i) = \varphi(T v_i) = 0$ say that $T v_i$ has no $u$-component --- i.e. that $\myspan{v_2, \ddd, v_n}$ is itself $T$-invariant. Nothing in the hypotheses provides that, and choosing a different complement does not help.

  The equivalence chain above avoids the issue by never naming an eigenvector at all. It works entirely with the \emph{failure} of injectivity and surjectivity, using that $\dual{S}$ is surjective exactly when $S$ is injective, and that in finite dimensions an operator is invertible as soon as it is injective. Existence statements are often easier to prove by ruling out invertibility than by exhibiting a witness.

  One caution on the chain: it uses finite-dimensionality twice, once for $\dual{(T-\lambda I)} = \dual{T} - \lambda I$ to behave and once for \qt{not surjective $\iff$ not invertible}. In infinite dimensions the last step fails.
\end{ainote}

\exercise{21}
\begin{xrcs}
  Suppose $T \in \linmap(V)$ is invertible.
  \begin{enumerate}
    \item Suppose $\lambda \in \myF$ with $\lambda \neq 0$. Prove that $\lambda$ is an eigenvalue of $T$ $\iff$ $\frac{1}{\lambda}$ is an eigenvalue of $T^{-1}$.
    \item Prove that $T$ and $T^{-1}$ have the same eigenvectors.
  \end{enumerate}
  \begin{xsol}
    Let $v \in V$ be such that $T(v) = \lambda v$. This is the case $\iff T^{-1} T(v) = T^{-1} (\lambda v)$ $\iff$ $T^{-1} (v) = \frac{1}{\lambda} v$. Note that in a field $\myF$, we can ``replace'' $\frac{1}{\lambda}\in \myF$ with its multiplicative inverse $\alpha \in \myF$. Since $\left(T^{-1}\right)^{-1} = T$, we have proven (a) and (b).
  \end{xsol}
\end{xrcs}

\exercise{22}
\begin{xrcs}
  Suppose $T \in \linmap(V)$ and there exist nonzero vectors $u, w \in V$ such that
  \begin{equation}
    Tu = 3w \quad \text{ and } \quad Tw = 3u.
  \end{equation}

  Prove that $3$ or $-3$ is an eigenvalue of $T$.

  \begin{xsol}
    Compute $T(u+w) = Tu + Tw = 3w + 3u = 3(u+w)$, and $T(u-w) = Tu - Tw = 3w - 3u = -3(u-w)$.
    Since $u \neq 0$, $(u+w) + (u-w) = 2u \neq 0$, so $u+w$ and $u-w$ cannot both be zero.
    If $u+w \neq 0$, then $u+w$ is an eigenvector with eigenvalue $3$.
    If $u-w \neq 0$, then $u-w$ is an eigenvector with eigenvalue $-3$.
    Hence, $3$ or $-3$ is an eigenvalue of $T$.
  \end{xsol}
\end{xrcs}

\begin{ainote}[Two things to get right here]
  \emph{You may not adjust $T$.} A tempting move when the case $u + w = 0$ is inconvenient is to redefine $T$, e.g. by declaring $Tu = 3(-w)$. That proves a statement about a \emph{different} operator. $T$ is given by the hypotheses $Tu = 3w$ and $Tw = 3u$, and the conclusion has to be reached for that $T$ --- the only freedom is in which \emph{vectors} we feed to it, which is why $u+w$ and $u-w$ are the right things to try.

  \emph{The exercise says \qt{or}, and it has to.} Neither eigenvalue can be forced on its own, because each degenerate case genuinely occurs. If $w = -u$ then $u + w = 0$ and only $-3$ is obtained; if $w = u$ then $u - w = 0$ and only $3$ is. Both are consistent with the hypotheses, e.g. on $V = \myF$ with $T$ multiplication by $3$ (take $u = w = 1$) or by $-3$ (take $u = 1, w = -1$). So the proof must split into cases, and the step that makes the split safe is $(u+w) + (u-w) = 2u \neq 0$: the two candidates cannot vanish simultaneously.

  Incidentally, $T^2 u = T(3w) = 9u$ and likewise $T^2 w = 9w$, so $T^2$ acts as $9I$ on $\myspan{u,w}$. That is the structural reason $\pm 3$ are the only candidates: the eigenvalues in play must square to $9$.
\end{ainote}

\exercise{23}
\begin{xrcs}
  Suppose $V$ is finite-dimensional and $T \in \linmap(V)$. Prove that $T$ has an eigenvalue $\iff$ there exists a subspace of $V$ of dimension $\mydim V - 1$ that is invariant under $T$.

  \begin{xprf}
    \Rightarrowdirection Suppose $T$ has an eigenvalue $\lambda$. By Exercise 15, $\lambda$ is an eigenvalue of $\dual{T} \in \linmap(\dual{V})$.
    Let $\varphi \in \dual{V} \setminus \{0\}$ be an eigenvector of $\dual{T}$ with eigenvalue $\lambda$, so $\dual{T}(\varphi) = \lambda \varphi$.
    Define $U = \mynull \varphi$. Since $\varphi \neq 0$, $\dim U = \dim V - 1$.
    For any $u \in U$, $\varphi(Tu) = (\dual{T}\varphi)(u) = (\lambda \varphi)(u) = \lambda \cdot 0 = 0$, so $Tu \in U$.
    Thus $U$ is an invariant subspace of dimension $\dim V - 1$.

    \Leftarrowdirection Suppose $U$ is a $T$-invariant subspace of $V$ with $\dim U = \dim V - 1$.
    The quotient space $V/U$ has dimension $\dim V - \dim U = 1$.
    The induced quotient operator $T/U \in \linmap(V/U)$ on a $1$-dimensional space is multiplication by some scalar $\lambda \in \myF$, by Exercise 3A.7.

    It remains to turn that $\lambda$ into an eigenvalue of $T$ itself, which does not follow merely from the words. Saying that $T/U$ is multiplication by $\lambda$ means
    \begin{equation}
      (T/U)(v + U) = Tv + U = \lambda v + U \quad \forall v \in V,
    \end{equation}
    i.e. $Tv - \lambda v \in U$ for every $v \in V$. In other words
    \begin{equation}
      \myrange (T - \lambda I) \subseteq U.
    \end{equation}

    Hence $\dim \myrange (T-\lambda I) \leq \dim U = \dim V - 1 < \dim V$, so $T - \lambda I$ is not surjective. Since $V$ is finite-dimensional, an operator on $V$ is surjective if and only if it is injective by \ref{thm: injectivity is equivalent to surjectivity}, so $T - \lambda I$ is not injective either. Therefore $\mynull (T - \lambda I) \neq \{0\}$, which is precisely to say that $\lambda$ is an eigenvalue of $T$.
  \end{xprf}
\end{xrcs}

\begin{ainote}[Why the obvious approach to $\Rightarrow$ does not work]
  The natural first attempt at the $\Rightarrow$ direction is to build the codimension-one subspace by hand: write $V = U \oplus \myspan{v_0}$ for some hyperplane $U$ and try to adjust a basis $u_1, \ddd, u_{n-1}$ of $U$ until $U$ becomes invariant. It does not come out. The obstruction is that each $T u_k$ decomposes as a part in $U$ \emph{plus} a part along $v_0$, and killing all $n-1$ of those $v_0$-components simultaneously is exactly as hard as the original problem --- there is no reason a basis change of $U$ should clear them.

  Passing to the dual sidesteps this entirely. A hyperplane is the same thing as the null space of a nonzero functional, so instead of choosing $n-1$ vectors one chooses a \emph{single} $\varphi \in \dual{V}$, and the invariance of $\mynull \varphi$ is the one condition $\dual{T} \varphi = \lambda \varphi$. The problem shrinks from adjusting a basis to finding one eigenvector --- of $\dual{T}$ rather than of $T$, which Exercise 15 supplies.

  This is the recurring payoff of duality: conditions on a subspace of codimension $k$ become conditions on a $k$-dimensional space of functionals, and small numbers are easier to handle than large ones.
\end{ainote}

\exercise{24}
\begin{xrcs}
  Suppose $T \in \linmap(V)$ and $\lambda_1, \ddd, \lambda_m$ are distinct eigenvalues of $T$. Prove that the sum of the corresponding eigenspaces is a direct sum:
  \begin{equation}
    E(\lambda_1, T) + \cdots + E(\lambda_m, T) = E(\lambda_1, T) \oplus \cdots \oplus E(\lambda_m, T).
  \end{equation}
  \begin{xprf}
    To show that the sum is direct, suppose $u_1 + \dots + u_m = 0$, where $u_j \in E(\lambda_j, T)$ for each $j \in \{1, \dots, m\}$.
    By \ref{thm: condition for a direct sum} it suffices to show that every $u_j = 0$.
    Suppose for contradiction that not all $u_j$ are zero.
    Discarding the terms where $u_j = 0$, we obtain a non-empty sum of nonzero eigenvectors corresponding to distinct eigenvalues that equals zero.
    However, by Theorem \ref{thm: linearly independent eigenvectors}, eigenvectors corresponding to distinct eigenvalues are linearly independent.
    Thus no non-trivial linear combination of them can equal zero, a contradiction.
    Therefore, $u_1 = \dots = u_m = 0$, which proves that the sum is a direct sum.
  \end{xprf}
\end{xrcs}

\begin{ainote}[The discarding step is what makes this work]
  It is worth being precise about why the zero terms are thrown away. \ref{thm: linearly independent eigenvectors} says that eigenvectors belonging to distinct eigenvalues are linearly independent, and \qt{eigenvector} means \emph{nonzero} by definition. So the theorem cannot be applied to the list $u_1, \ddd, u_m$ as it stands: some $u_j$ may be $0$, and $0$ is not an eigenvector.

  Discarding them repairs exactly that. What is left is a genuine list of eigenvectors for distinct eigenvalues summing to $0$ with all coefficients equal to $1 \neq 0$, which contradicts their independence. If \emph{every} term was discarded there is nothing to contradict --- but that is the case $u_1 = \cdots = u_m = 0$, which is what we wanted anyway.
\end{ainote}

\exercise{25}
\begin{xrcs}
  Suppose $T \in \linmap (V)$ and $u,w$ are eigenvectors of $T$ such that $u + w$ is also an eigenvector of $T$. Prove that $u$ and $w$ are eigenvectors of $T$ corresponding to the same eigenvalue.

  \begin{xprf}
    Let $u,w \in V \setminus \{0\}$ such that $T u = \lambda_u u$, $T w = \lambda_w w$, and $T (u+w) = \lambda_{u+w} (u+w)$. Since
    \begin{equation}
      T (u + w) =  T u + T w = \lambda_u u + \lambda_w w = \lambda_{u+w} (u + w) = \lambda_{u+w} u + \lambda_{u+w} w,
    \end{equation}

    it follows that
    \begin{equation}
      \left(\lambda_u -  \lambda_{u+w}\right) u + \left(\lambda_w -  \lambda_{u+w}\right) w = 0.
    \end{equation}

    If $u$ and $w$ are linearly independent, the coefficients must vanish, so $\lambda_u =  \lambda_{u+w} = \lambda_w$. Otherwise, if $u$ and $w$ are linearly dependent, we can write $u = \alpha w$ for some $\alpha \neq 0$. Applying $T$ to $u$ in two ways gives
    \begin{equation}
      \begin{aligned}
        T u &= T(\alpha w) = \alpha T w = \alpha \lambda_w w \text{ and }  \\
        T u &= \lambda_u u= \lambda_u \alpha w.
      \end{aligned}
    \end{equation}

    Hence, $\lambda_u = \lambda_w$.
  \end{xprf}
\end{xrcs}

\exercise{26}
\begin{xrcs}
  Suppose $T \in \linmap(V)$ and $u, v \in V$ are nonzero eigenvectors of $T$ corresponding to distinct eigenvalues $\lambda \neq \mu$. Prove that $u+v$ is not an eigenvector of $T$.

  \begin{xprf}
    Suppose for the sake of contradiction that $u+v$ is an eigenvector of $T$.
    By Exercise 25, $u$ and $v$ must be eigenvectors of $T$ corresponding to the same eigenvalue.
    However, $u$ corresponds to $\lambda$ and $v$ corresponds to $\mu$ with $\lambda \neq \mu$, a contradiction.
    Therefore, $u+v$ cannot be an eigenvector of $T$.
  \end{xprf}
\end{xrcs}

\exercise{33}
\begin{xrcs}
  Suppose $T \in \linmap(V)$ and $m$ is a positive integer.
  \begin{enumerate}
    \item Prove that $T$ is injective $\iff$ $T^m$ is injective.
    \item Prove that $T$ is surjective $\iff$ $T^m$ is surjective.
  \end{enumerate}

  \begin{xprf}
    We show (a) and (b) both with induction, first in the $\Rightarrow$ direction, and then in the $\Leftarrow$ direction.
    \begin{enumerate}
      \item {
        \Rightarrowdirection Suppose $T$ is injective. For $m=1$ the statement is trivial ($T^1 = T$). Fix  $m \geq 1$ and suppose as induction hypothesis that $T^m$ is injective. If $T^{m+1}(v_1) = T^{m+1} (v_2)$, then
        \begin{equation}
          T(T^{m}(v_1)) = T(T^{m} (v_2)).
        \end{equation}

        Since $T$ is injective, $T^{m}(v_1) = T^{m} (v_2)$. The induction hypothesis forces $v_1 = v_2$. Therefore, $T^{m+1}$ is injective as well. (In general, function composition is injective. If $f: X \to Y$ and $g: Y \to Z$ are injective, so is $g \circ f : X \to Z$. So, a finite composition of injective maps is injective.)

        \Leftarrowdirection Assume $T^m$ is injective for some fixed $m \geq 1$. If $m=1$ we are done. If $m>1$, take $v_1, v_2 \in V$ with $T(v_1) = T(v_2)$. Apply $T^{m-1}$ to both sides to get $T^{m-1} (T(v_1)) = T^{m-1} (T(v_2))$. Injectivity of $T^{m}$ forces $v_1 = v_2$. Thus, $T$ is injective.
      }
      \item{
        \Rightarrowdirection Suppose $T$ is surjective. For $m=1$ the statement is trivial. Fix $m \geq 1$ and assume that $T^m$ is surjective. Let $y\in V$ be arbitrary. Since $T$ is surjective, there exists $z\in V$ with $T(z)=y$. By the induction hypothesis, there exists $v\in V$ with $T^m(v)=z$. Thus,
        \begin{equation}
          T^{m+1}(v)=T(T^m(v))=T(z)=y.
        \end{equation}

        Hence, $T^{m+1}$ is surjective.

        (Or: A finite composition of surjective maps (not necessarily linear) is surjective.)

        \Leftarrowdirection Assume $T^m$ is surjective for some fixed $m\ge 1$. Let $y\in V$ be arbitrary. Choose $v\in V$ with $T^m(v)=y$ (always possible, since $T^m$ is surjective). Put $x:=T^{m-1}(v)$. Then
        \begin{equation}
          y=T(x)=T(T^{m-1}(v))=T^m(v).
        \end{equation}

        Hence, $T$ is surjective.
      }
    \end{enumerate}
  \end{xprf}
\end{xrcs}

\exercise{40}
\begin{xrcs}
  Suppose $S,T \in \linmap (V)$, and $S$ is invertible. Suppose $p \in \polyn(\myF)$ is a polynomial. Prove that
  \begin{equation}
    p(STS^{-1}) = S p(T) S ^{-1}.
  \end{equation}

  \begin{xprf}
    We first want to prove that $(STS^{-1})^m = ST^m S^{-1}$ for every $m \in \nat_{\geq 0}$. To this end, we use induction over $m$. For $m=0$, the statement holds because
    \begin{equation}
      (STS^{-1})^0 = I = SIS^{-1} = ST^0S^{-1}.
    \end{equation}

    If $m=1$ the claim is immediate. Hence, assume $m \geq 1$. Then
    \begin{equation}
      \begin{aligned}
        (STS^{-1})^{m+1} &= STS^{-1} (STS^{-1}) ^m \\
          &= ST\underbrace{S^{-1} S}_{I}T^m S^{-1} \\
          &= S T^{m+1} S^{-1},
      \end{aligned}
    \end{equation}

    This completes the induction. Let $p(z)=a_0 z^0 + a_1 z^1 + \cdots + a_m z^m, \; a_j \in \myF$. Then
    \begin{equation}
      \begin{aligned}
        p(STS^{-1})
          &= a_0 I \phantom{ A A A A A A  } + a_1(STS^{-1}) + a_2 (STS^{-1})^2 + \cdots + a_m (STS^{-1})^m
            & = \textstyle \sum_{j=0}^m a_j\,(STS^{-1})^j \\
          &= a_0 I \underbrace{S S^{-1} }_{= \, I \, = (SIS^{-1})^0} + \, a_1(STS^{-1}) + a_2 (ST^2S^{-1}) + \cdots + a_m (ST^mS^{-1})
            & = \textstyle \sum_{j=0}^m a_j\,(ST^jS^{-1}) \\
          &= S (a_0 I)  S^{-1} + S(a_1T)S^{-1} + S(a_2 T^2)S^{-1} + \cdots + S(a_mT^m)S^{-1}
            & = \textstyle \sum_{j=0}^m S(a_jT^j)S^{-1}\\
          &= S \left( a_0 I  + a_1T + a_2 T^2 + \cdots + a_mT^m \right)S^{-1}
            & = S \left(\textstyle \sum_{j=0}^m a_jT^j\right) S^{-1} \\
          &= S p(T) S^{-1}.
      \end{aligned}
    \end{equation}

    Here we have been making use of the distributive property of linear maps, see  \ref{thm: algebraic properties of products of linear maps}. And we used that scalar multiples of the identity (i.e. $a_0 I$) commute with every operator.
  \end{xprf}
\end{xrcs}

